% !TEX root = ../../../main.tex
% 来源：中学数学实验教材·第二册（几何）
% 章节：实验几何 / 基本作图
% 原始文件：1.tex，行 3955-4015
% 类型：tikz
% ---
\begin{tikzpicture}[scale=.7]
\begin{scope}
\draw(0,0) rectangle (2,2);
\draw(2,0) rectangle (4,2);
\draw(4,0) rectangle (6,2);

\foreach \x in {0,2,4}
{
	\draw(\x+1,0)--(\x+1,2);
	\draw(\x,0) arc (-90:90:1);
	\draw(\x,0) arc (180:0:1);
	\draw(\x,2) arc (90:-90:1);
	\draw(\x,2) arc (180:360:1);
    \draw(\x+2,0) arc (270:90:1);
}
\draw(0,1)--(6,1);

\end{scope}

\begin{scope}[xshift=10cm]
\tkzDefPoints{0/0/A, 2/0/B,1/1.732/C}
\foreach \x in {A,B,C}
{
	\draw(\x) circle (1);
}
\tkzDrawPolygon[fill=white](A,B,C)
\end{scope}
\begin{scope}[xshift=2cm, yshift=-3cm]
	\draw(0,0)circle(2);
\draw(30:2) arc (90:210:2);
\draw(30:2) arc (-30:-150:2);
\draw(-30:2) arc (-90:-210:2);
\draw(-30:2) arc (30:150:2);
\draw(90:2) arc (30:-90:2);
\draw(-90:2) arc (-30:90:2);

\end{scope}

\begin{scope}[xshift=6cm, yshift=-3cm]
\draw(2,0)circle(2);
\foreach \x in {1,2,3}
{
	\draw(0,0) arc (180:0:\x/2);
	\draw(4,0) arc (0:-180:\x/2);
}
\draw(0,0)--(4,0);
\end{scope}

\begin{scope}[xshift=14cm, yshift=-3cm]
	\draw(0,0)circle(2);
\draw(-2,0)--(2,0);
\draw(0,-2)--(0,2);
\foreach \x in {0,2}
{
	\draw(\x-2,0) arc (-180:0:.5);
	\draw(\x-1,0) arc (180:0:.5);
	\draw(0,\x) arc (90:270:.5);
	\draw(0,\x-1) arc (90:-90:.5);
}
\end{scope}
\end{tikzpicture}
